
% =============================================================
\section{Physikalische Grundlagen}
% =============================================================

Der Compton-Effekt beschreibt die Streuung eines Photons an einem zunächst ruhenden Elektron. 
Um die beobachtete Wellenlängenänderung zu verstehen, müssen sowohl die Energie 
als auch der Impuls des Photons berücksichtigt werden.

% -------------------------------------------------------------
\subsection{Energie und Impuls des Photons}
% -------------------------------------------------------------
Ein Photon besitzt nach der Quantentheorie die Energie
\[
E = h\nu
\]
und den Impuls
\[
p = \frac{E}{c} = \frac{h\nu}{c} = \frac{h}{\lambda}.
\]
Obwohl Photonen keine Ruhemasse besitzen, tragen sie einen Impuls, 
der bei Streuprozessen übertragen werden kann. 

\begin{DidaktikBox}[Energie im Gegensatzt zur klassischen Wellenbeschreibung]
	In der klassischen Wellentheorie wird Energie über die Amplitude transportiert, 
	nicht über einzelne Quanten. 
	Der Compton-Effekt zeigt, dass Licht in einzelnen Energiepaketen \( h\nu \) auftritt 
	– ein klarer Hinweis auf den Teilchencharakter des Lichts.
\end{DidaktikBox}

% -------------------------------------------------------------
\subsection{Stoß zwischen Photon und Elektron}
% -------------------------------------------------------------
Ein Photon trifft auf ein ruhendes Elektron (Masse \( m_e \)). 
Nach der Streuung besitzt das Photon eine geringere Energie \( h\nu' \) 
und wird unter dem Winkel \( \theta \) gestreut. 
Das Elektron erhält einen Rückstoßimpuls \( \vec{p}_e \).

\begin{PhysikBox}[Streugeometrie]
	\small
	Die Streugeometrie kann vektoriell beschrieben werden durch
	\[
	\vec{p}_\gamma = \vec{p}_{\gamma'} + \vec{p}_e.
	\]
	In Komponenten:
	\[
	p_\gamma = p_{\gamma'}\cos\theta + p_{e,x}, \qquad
	0 = p_{\gamma'}\sin\theta - p_{e,y}.
	\]
\end{PhysikBox}

\begin{HinweisBox}[Nichtelastischer Stoß]
	Obwohl der Stoß Energie- und Impulserhaltung folgt, 
	ist er \emph{nicht elastisch} im klassischen Sinn, 
	da ein Teil der Photonenenergie in kinetische Energie des Elektrons übergeht.
\end{HinweisBox}

% -------------------------------------------------------------
\subsection{Energie- und Impulserhaltung kombiniert}
% -------------------------------------------------------------
Nach der Energieerhaltung gilt:
\[
h\nu + m_e c^2 = h\nu' + \sqrt{(p_e c)^2 + (m_e c^2)^2}.
\]
Kombiniert man dies mit den Impulsbeziehungen, 
erhält man nach einigen Umformungen die bekannte \emph{Compton-Formel}:
\[
\lambda' - \lambda = \frac{h}{m_e c}\,(1 - \cos\theta).
\]

\begin{MatheBox}[Die Compton-Wellenlänge]
	Die Größe
	\[
	\lambda_C = \frac{h}{m_e c} \approx \SI{2.43e-12}{\meter}
	\]
	wird als \textbf{Compton-Wellenlänge} bezeichnet 
	und charakterisiert die Skala, auf der der Effekt messbar ist.
\end{MatheBox}

% -------------------------------------------------------------
\subsection{Physikalische Interpretation}
% -------------------------------------------------------------
Die Wellenlängenänderung hängt nur vom Streuwinkel \( \theta \) ab 
und ist maximal bei Rückstreuung (\( \theta = 180^\circ \)). 
Dies bestätigt, dass Photonen Energie und Impuls wie Teilchen übertragen, 
obwohl sie sich gleichzeitig wellenartig ausbreiten.

\begin{DidaktikBox}[Was der Compton-Effekt zeigt]
	Der Compton-Effekt beweist, dass Licht nicht allein als Welle verstanden werden kann. 
	Er verlangt eine Beschreibung durch \emph{Teilchen mit Energie und Impuls} – 
	die Photonen. Damit wurde die klassische Elektrodynamik entscheidend erweitert.
\end{DidaktikBox}
